% Generated draft: requires mathematical and source review.
\documentclass[11pt]{article}
\usepackage[margin=25mm]{geometry}
\usepackage{fontspec}
\setmainfont{Latin Modern Roman}
\usepackage{amsmath,amssymb,mathtools,unicode-math}
\setmathfont{Latin Modern Math}
\usepackage{enumitem,needspace,xurl,hyperref}
\setlength{\parindent}{0pt}
\setlength{\parskip}{5pt}
\setlength{\emergencystretch}{3em}
\newcommand{\N}{\mathbb{N}}
\newcommand{\Z}{\mathbb{Z}}
\newcommand{\Q}{\mathbb{Q}}
\newcommand{\R}{\mathbb{R}}
\newcommand{\C}{\mathbb{C}}
\newcommand{\F}{\mathbb{F}}
\newcommand{\A}{\mathbb{A}}
\newcommand{\PP}{\mathbb{P}}
\newcommand{\E}{\mathbb{E}}
\newcommand{\eps}{\varepsilon}
\newcommand{\dd}{\,\mathrm{d}}
\begin{document}
\section*{UnsolvedMath ID: 30003126}
\textbf{Source problem code:} catalogue entry
\subsection{Definitions and mathematical statement}
The record leaves the curve model, the set \(S\), and the convention for integrality implicit. Under the usual interpretation, let \(k\) be a number field, let \(C/k\) be a smooth, projective, geometrically integral curve of genus \(2\), let \(\phi\in k(C)\setminus k\), and let \(S\) be a finite set of places of \(k\) containing all archimedean places. These conventions make the intended question precise but require confirmation against the original source.

Define the ring of \(S\)-integers by
\[
\mathcal O_{k,S}=\{a\in k: v(a)\geq 0\text{ for every nonarchimedean place }v\notin S\}.
\]
Let \(D_\infty\) be the pole divisor of \(\phi\), and interpret the points that are \(S\)-integral with respect to \(\phi\) as
\[
I(C,\phi,k,S)=\{P\in C(k)\setminus\operatorname{Supp}(D_\infty):\phi(P)\in\mathcal O_{k,S}\}.
\]
In particular, poles of \(\phi\) are excluded, and the points considered are \(k\)-rational points.

\textbf{Effectivity question.} Is there a single algorithm which, for every such input \((k,C,\phi,S)\) given by effective finite descriptions, terminates and outputs exactly the set \(I(C,\phi,k,S)\)? The source also asks for an explicit bound for all these points. A precise height-bound version asks whether, given additionally an effectively specified projective embedding \(\iota:C\hookrightarrow\PP^N_k\), one can compute a finite bound \(B\) from the input such that
\[
 h_{\iota}(P)\leq B\qquad\text{for every }P\in I(C,\phi,k,S),
\]
where \(h_{\iota}\) is the absolute logarithmic projective Weil height induced by \(\iota\). The source does not specify a height, an input representation, or a required form of the bound; this height formulation is an explicit interpretation of its alternative request, not a separately supplied formulation.
\subsection{Short English statement}
For every genus-two curve over a number field, nonconstant rational function, and finite set of allowed places, ask for an effective procedure to list all rational points where the function is finite and takes an \(S\)-integer value, or for a computable height bound on those points. The curve-model and integrality conventions above require source confirmation. The supplied record labels the problem open; this draft does not independently verify that status.
\subsection{Sources}
\begin{itemize}
\item Local UnsolvedMath registry, permanent ID \(30003126\), source code \texttt{catalogue entry}, title \emph{Effective Siegel Theorems for Genus-Two Curves}, in the Oberwolfach Reports Open Problems collection. No individual proposer is supplied. The record gives proposed year \(2016\).
\item Supplied report citation: \emph{Diophantische Approximationen (2017)}, Oberwolfach Reports, pp.\ 1099--1187, DOI \texttt{10.4171/owr/2016/21}. \url{https://doi.org/10.4171/owr/2016/21}. The supplied background describes these reports as documenting mathematical workshops and their open-problem discussions. No precise problem passage is supplied.
\item Levin, \emph{Siegel's Theorem and the Shafarevich Conjecture}. Supplied literature link: \url{https://arxiv.org/abs/1109.6070}.
\item Bugeaud, Mignotte, Siksek, Stoll, and Tengely, \emph{Integral Points on Hyperelliptic Curves}. Supplied literature link: \url{https://arxiv.org/abs/0801.4459}.
\item Literature assessment checked at \texttt{2026-08-21}: the supplied curation labels the problem open and states that no general effective Siegel theorem for genus at least two is known, that special genus-two families can be handled effectively, and that the stated general effectivity problem remains unresolved. This is supplied dated curation derived from the OWR collection, not an independently verified literature assessment; its status and cited claims require independent verification.
\item The registry assigns statement status \texttt{corrected\_verified} and says that mathematical content was preserved while notation, typography, or omitted definitions were repaired. This is a supplied editorial qualification, not independent mathematical verification of this draft.
\end{itemize}
\end{document}
